\documentclass[11pt]{article}
\usepackage{series}
\usepackage{booktabs}
\usepackage{array}

\title{\textbf{ETH and Many-Body Localization Through an Operational Basin Diagnostic}\\
\large A Model-Dependent Bridge with Exact Coarse-Dynamics Identities}
\author{Peter Nero}
\date{Version 2, July 2026}

\hypersetup{
  pdftitle={ETH and Many-Body Localization Through an Operational Basin Diagnostic},
  pdfauthor={Peter Nero}
}

\DeclareMathOperator{\Tr}{Tr}
\newcommand{\Dens}{\mathcal D}
\newcommand{\Obs}{\mathcal O}
\newcommand{\Micro}{\mathcal M}
\newcommand{\Eff}{\mathcal K}
\newcommand{\EE}{\mathcal E}
\newcommand{\II}{\mathcal I}
\newcommand{\Dyn}{\mathcal A}
\newcommand{\TT}{\mathcal T}
\newcommand{\id}{\mathbf 1}
\newcommand{\rhoMC}{\rho_{\mathrm{mc}}}

\begin{document}
\maketitle

\begin{abstract}
This paper replaces a proposed universal equivalence among eigenstate
thermalization, many-body localization, and projection-induced basins by
a model-dependent operational bridge.  We fix a finite random-field XXZ
chain, an energy window, an initial-state class, and a family of local
observables.  Standard diagnostics---eigenstate matrix elements,
adjacent-gap ratios, imbalance, and half-chain entanglement---are kept
distinct.  For a coarse channel followed by microscopic unitary
evolution, we derive the exact failure of the reduced maps to form a
semigroup.  We prove that microscopic unitary dynamics preserves trace
distance, that uniform contraction of an effective map is a sufficient
but not necessary thermalization criterion, and that the mean-square
thermalization error separates exactly into diagonal bias and temporal
fluctuation.  For a quasi-local integral of motion whose commutator with
the Hamiltonian is bounded by $\eta$, we derive a memory bound linear in
$\eta |t|$.  These results explain when ``mixing basin'' and
``memory basin'' are useful operational descriptions without identifying
them with ETH or MBL by definition.  No universal logistic or
Kramers-type knee follows for a closed quantum chain; a crossover must be
estimated with a declared finite-size scaling model and uncertainty
analysis.  The Modal Triplet Theory interpretation is therefore a
conditional pullback problem: its coherent projector, Hessian gap, and
capacity margin become relevant only after an explicit intertwiner maps
them to the Hamiltonian, coarse channel, observables, and diagnostics
used here.
\end{abstract}

\section*{Version 2 Revision Note}

\begin{description}
\item[Supersedes.] Version 1, \emph{ETH and Many-Body Localization as a
Single Shadow-Bridge Problem}.
\item[Reason.] The earlier version treated thermalization, localization,
and monitored transitions as universal basin equivalences; it did not
declare a Hamiltonian family; and it imposed a Kramers-type crossover
without deriving the stochastic escape process required for that law.
\item[Resolution.] Version 2 fixes a random-field XXZ family and separate
diagnostics, replaces equivalences by exact identities and sufficient
criteria, derives a quasi-LIOM memory estimate, and gives a finite-size
comparison protocol with no universal knee assumption.
\item[Retained content.] The basin language is retained as a useful
operational summary of coarse mixing and memory when its metric, map,
time window, and state class are declared.
\item[Open boundary.] No current theorem derives the XXZ Hamiltonian,
coarse channel, disorder law, thermal reference, or scaling function from
one selected MTT source.  The thermodynamic ETH--MBL boundary and a
general asymptotic MBL phase are not established here.
\end{description}

\section{Introduction}

Thermalization and many-body localization are not merely opposite labels.
They are statements about different mathematical objects.  Eigenstate
thermalization concerns matrix elements of chosen observables in energy
eigenstates \cite{deutsch1991,srednicki1994}.  Dynamical equilibration
concerns the evolution of selected initial states and observables
\cite{reimann2008}.  Many-body localization may be diagnosed through
level statistics, eigenstate entanglement, transport, persistent local
memory, or quasi-local integrals of motion
\cite{oganesyan2007,pal2010,imbrie2016}.  These diagnostics are related
under additional hypotheses, but they are not definitions of one another.

The earlier version of this paper compressed all of them into a single
picture: one dominant basin represented ETH, while many fragmented basins
represented MBL.  The picture is suggestive but too strong.  A basin is
defined only after a state space, dynamics, metric, resolution, and time
window have been chosen.  Different choices can produce different
partitions of the same microscopic system.  Integrable systems can retain
memory without being MBL.  Conversely, finite systems can show long
crossovers without establishing a thermodynamic phase.

The revised aim is narrower and more useful.  We ask when basin language
is a faithful \emph{operational summary} of independently specified
many-body diagnostics.  The answer has three parts:
\begin{enumerate}
\item a declared Hamiltonian and diagnostic record;
\item exact statements about the relation between unitary and coarse
dynamics; and
\item an explicit interface that MTT would have to satisfy before its
projector or stability margin can explain the observed many-body regime.
\end{enumerate}

\section{Claim tier and scope}

\subsection{What is proved}

All exact results in this paper are finite-dimensional.  They hold for a
fixed chain length, a fixed disorder realization unless averaging is
explicitly stated, and declared maps and observables.  We prove:
\begin{enumerate}
\item the exact semigroup defect of an evolve--coarse-grain construction;
\item preservation of trace distance by microscopic unitary evolution;
\item exponential convergence under an assumed effective contraction;
\item an exact decomposition of mean-square thermalization error; and
\item a memory estimate for an operator that nearly commutes with the
Hamiltonian.
\end{enumerate}

\subsection{What is not proved}

We do not prove ETH or MBL for the entire XXZ family.  We do not infer an
asymptotic transition from finite exact diagonalization, derive a
universal critical disorder, or identify a Hamiltonian level-spacing gap
with an MTT fixed-point Hessian gap.  Rigorous localization results exist
for restricted one-dimensional spin-chain settings and assumptions
\cite{imbrie2016}; rare-region and avalanche mechanisms make broader
stability questions model- and dimension-dependent
\cite{deroeck2017,morningstar2022}.  The present bridge respects that
scope.

\section{A complete many-body problem record}

\subsection{Hamiltonian family}

For $L\geq 2$, let
\[
 \Micro_L=(\bbC^2)^{\otimes L}
\]
and impose open boundaries.  We use the random-field XXZ Hamiltonian
\begin{equation}\label{eq:xxz}
 H_L(\omega;W)
 =
 J\sum_{j=1}^{L-1}
 \left(
 S_j^xS_{j+1}^x+S_j^yS_{j+1}^y+
 \Delta S_j^zS_{j+1}^z
 \right)
 +\sum_{j=1}^{L}h_j(\omega)S_j^z ,
\end{equation}
where $J\neq0$, $\Delta\in\bbR$, and the $h_j$ are independent uniform
random variables on $[-W,W]$.  Numerical comparisons must fix the
symmetry sector, here total $S^z=0$ for even $L$, because unresolved
symmetry blocks can create misleading degeneracies and gap statistics.
They must also state the boundary condition, energy-density window,
disorder seed or sample list, number of samples, and diagonalization
tolerance.

\subsection{Diagnostic record}

Let $H_L|n\rangle=E_n|n\rangle$ within the selected sector and energy
window.  A minimally complete record contains the following distinct
rows.

\paragraph{Eigenstate matrix elements.}
For each declared local observable $O$,
\begin{equation}\label{eq:eth-matrix}
 O_{mn}=\langle m|O|n\rangle .
\end{equation}
An ETH analysis specifies a smooth microcanonical target for $O_{nn}$,
the width of diagonal residuals, and a scaling test for off-diagonal
elements.  A statement that trajectories approach one coarse state is
not by itself the ETH ansatz.

\paragraph{Adjacent-gap ratio.}
For ordered levels, define $\delta_n=E_{n+1}-E_n$ and
\begin{equation}\label{eq:ratio}
 r_n=\frac{\min(\delta_n,\delta_{n+1})}
 {\max(\delta_n,\delta_{n+1})}.
\end{equation}
The familiar reference means are approximately $0.5307$ for the
Gaussian orthogonal ensemble and $2\log 2-1\approx0.3863$ for Poisson
statistics \cite{oganesyan2007,pal2010}.  These are ensemble diagnostics,
not proofs of ETH or MBL.

\paragraph{Memory and entanglement.}
For the N\'eel state $|\psi_{\mathrm N}\rangle$, one may use
\begin{equation}\label{eq:imbalance}
 I_L(t)=\frac{2}{L}\sum_{j=1}^{L}(-1)^j
 \langle\psi_{\mathrm N}|S_j^z(t)|\psi_{\mathrm N}\rangle ,
\end{equation}
and the von Neumann entropy $S_{L/2}(t)$ of one half of the chain.
Persistent imbalance, slow entanglement growth, Poisson-like levels, and
area-law eigenstates are complementary evidence.  None should silently
stand in for all the others.

\begin{definition}[Comparison capsule]\label{def:capsule}
A finite-volume ETH--MBL comparison capsule is the tuple
\[
 \mathfrak C_L=(H_L,\mathcal S_L,\mathcal W_L,\mathcal R_L,
 \Obs_L,\mathcal I_L,\mathcal D_L,\mathcal U_L),
\]
where $\mathcal S_L$ is the symmetry sector, $\mathcal W_L$ the energy
window, $\mathcal R_L$ the disorder law and realized samples,
$\Obs_L$ the observable family, $\mathcal I_L$ the initial-state class,
$\mathcal D_L$ the diagnostics, and $\mathcal U_L$ the numerical and
statistical uncertainty record.
\end{definition}

Without these rows, a reported crossover cannot be reproduced or compared
across sizes.

\section{Coarse dynamics and its exact defect}

\subsection{Typed construction}

Let $\Dens(\Micro_L)$ and $\Dens(\Eff_L)$ denote microscopic and effective
density operators.  Let
\[
 \EE:\Dens(\Micro_L)\to\Dens(\Eff_L),\qquad
 \II:\Dens(\Eff_L)\to\Dens(\Micro_L)
\]
be completely positive trace-preserving maps with
$\EE\II=\id_{\Eff}$.  The microscopic channel is
\[
 \Dyn_t(X)=U_tXU_t^\dagger,\qquad U_t=e^{-itH_L},
\]
and the reduced family is
\begin{equation}\label{eq:reduced}
 \TT_t=\EE\Dyn_t\II .
\end{equation}
This is a family of quantum channels.  It is not automatically a
semigroup.

\begin{theorem}[Exact reduced-semigroup defect]\label{thm:defect}
For all $s,t$,
\begin{equation}\label{eq:defect}
 \TT_{t+s}-\TT_t\TT_s
 =
 \EE\Dyn_t(\id_{\Micro}-\II\EE)\Dyn_s\II .
\end{equation}
Consequently, $\{\TT_t\}$ is a semigroup exactly when the right-hand side
vanishes for every $s,t$ on the effective state span.
\end{theorem}

\begin{proof}
Using $\Dyn_{t+s}=\Dyn_t\Dyn_s$,
\[
\TT_{t+s}-\TT_t\TT_s
=
\EE\Dyn_t\Dyn_s\II-\EE\Dyn_t\II\EE\Dyn_s\II,
\]
and factoring the middle term gives \cref{eq:defect}.
\end{proof}

The defect has a direct meaning.  Evolving an embedded effective state
can generate microscopic information outside the selected representative
subspace.  Coarse-graining after the first interval discards that
information, while uninterrupted evolution can later return some of it.
A Markov approximation is therefore an additional approximation, not a
consequence of projection notation.

\begin{corollary}[Norm control of the defect]\label{cor:defect}
For any submultiplicative superoperator norm,
\[
 \norm{\TT_{t+s}-\TT_t\TT_s}
 \leq
 \norm{\EE}\,\norm{\Dyn_t}\,
 \norm{(\id_{\Micro}-\II\EE)\Dyn_s\II}.
\]
\end{corollary}

\subsection{Where contraction can enter}

Write $D(\rho,\sigma)=\frac12\norm{\rho-\sigma}_1$ for trace distance.

\begin{proposition}[Unitary evolution is not strict mixing]\label{prop:isometry}
For all microscopic states $\rho,\sigma$,
\[
 D(\Dyn_t\rho,\Dyn_t\sigma)=D(\rho,\sigma).
\]
If $\EE\II=\id_{\Eff}$ and both maps are channels, then
$D(\II\rho,\II\sigma)=D(\rho,\sigma)$.  Any strict reduction of
distinguishability under \cref{eq:reduced} therefore comes from the
coarse channel acting after microscopic excursion, not from unitary
evolution alone.
\end{proposition}

\begin{proof}
The trace norm is invariant under unitary conjugation.  Contractivity of
trace distance under channels gives
\[
D(\rho,\sigma)=D(\EE\II\rho,\EE\II\sigma)
\leq D(\II\rho,\II\sigma)\leq D(\rho,\sigma),
\]
so both inequalities are equalities.
\end{proof}

\begin{theorem}[Uniform effective contraction is sufficient]\label{thm:contraction}
Let $\TT$ be a sampled effective channel, let $\rho_\ast$ be a fixed
point, and let $\mathcal B\subset\Dens(\Eff_L)$ be invariant.  If
\[
 d(\TT\rho,\rho_\ast)\leq q\,d(\rho,\rho_\ast)
 \quad\text{for all }\rho\in\mathcal B,\qquad 0\leq q<1,
\]
then
\[
 d(\TT^n\rho,\rho_\ast)\leq q^n d(\rho,\rho_\ast).
\]
If $d$ controls every observable in $\Obs_L$, their expectation values
converge at the same rate.
\end{theorem}

\begin{proof}
Iteration gives the first inequality.  The observable statement follows
from the declared domination of expectation-value differences by $d$.
\end{proof}

This is a useful basin theorem, but it is one-way.  ETH does not require
one globally contractive channel, and a chosen coarse channel can be
contractive even when the microscopic eigenstate matrix elements fail an
ETH test.

\section{Thermalization without a basin tautology}

Fix an initial state $\rho$, a bounded observable $O$, and an energy
window with microcanonical state $\rhoMC$.  Suppose the infinite-time
Ces\`aro averages below exist and define the diagonal ensemble
\[
 \omega=\overline{\rho(t)}
 =\lim_{T\to\infty}\frac1T\int_0^T\rho(t)\,\dd t.
\]
Define
\[
 \epsilon_{\mathrm{diag}}(O)
 =\Tr[O(\omega-\rhoMC)]
\]
and the temporal fluctuation
\[
 f_O(t)=\Tr[O(\rho(t)-\omega)].
\]

\begin{theorem}[Exact thermalization-error decomposition]\label{thm:decomposition}
The mean-square deviation from the microcanonical value is
\begin{equation}\label{eq:decomposition}
 \overline{\left|\Tr[O(\rho(t)-\rhoMC)]\right|^2}
 =
 |\epsilon_{\mathrm{diag}}(O)|^2+
 \overline{|f_O(t)|^2}.
\end{equation}
\end{theorem}

\begin{proof}
The quantity inside the absolute value is
$\epsilon_{\mathrm{diag}}(O)+f_O(t)$.  By definition
$\overline{f_O}=0$, so the two cross terms vanish after time averaging.
\end{proof}

\Cref{eq:decomposition} separates two obligations often conflated by
basin language.  Small diagonal bias requires the populated eigenstates
to give the correct equilibrium expectation.  Small temporal variance
requires dephasing or an effective-dimension condition.  A useful
thermalization claim must control both, for declared states and
observables.  Existing equilibration theorems provide such bounds under
specific spectral and population hypotheses \cite{reimann2008}; the
basin metaphor does not replace them.

\begin{definition}[Operational thermal basin]\label{def:thermal-basin}
For tolerances $\varepsilon_{\rm d},\varepsilon_{\rm f}>0$, define
\[
\mathcal B_{\rm th}(O)=
\left\{\rho:
|\epsilon_{\rm diag}(O)|\leq\varepsilon_{\rm d},\
\overline{|f_O|^2}\leq\varepsilon_{\rm f}^2
\right\}.
\]
For a finite observable family, intersect these sets over $O\in\Obs_L$.
\end{definition}

This definition is deliberately diagnostic.  Membership records verified
thermal behavior at the selected resolution; it does not explain why the
Hamiltonian has that behavior.

\section{Memory and quasi-local integrals of motion}

Many-body localization is often described using quasi-local integrals of
motion.  The following elementary estimate makes the time scale explicit.

\begin{theorem}[Quasi-conserved memory bound]\label{thm:liom}
Let $\tau=\tau^\dagger$ be bounded and suppose
\[
 \norm{[H_L,\tau]}\leq\eta.
\]
Then, for every density operator $\rho$ and every $t\in\bbR$,
\begin{equation}\label{eq:liom}
 \left|\Tr[\rho\,\tau(t)]-\Tr[\rho\,\tau]\right|
 \leq \eta |t|.
\end{equation}
\end{theorem}

\begin{proof}
The Heisenberg equation gives
\[
\frac{\dd}{\dd t}\tau(t)=i\,U_t^\dagger[H_L,\tau]U_t.
\]
After integration, trace duality, and unitary norm invariance,
\[
\left|\Tr\!\left[\rho(\tau(t)-\tau)\right]\right|
\leq\int_0^{|t|}\norm{[H_L,\tau]}\,\dd s
\leq\eta|t|.
\]
\end{proof}

If $\tau$ is quasi-local near site $i$ and has a nonzero overlap with a
measured local spin, \cref{eq:liom} yields an operational memory lower
bound up to the approximation and overlap errors.  A family with
$\eta_L\to0$ can support growing memory times.  The converse does not
hold without further hypotheses: exact symmetries, integrability,
prethermal conservation, or kinetic bottlenecks can also preserve
memory.

\begin{definition}[Operational memory basin]\label{def:memory-basin}
For an initial-state class $\mathcal I_L$, observable $O$, threshold
$m_0>0$, and time window $[0,T_L]$, define
\[
\mathcal B_{\rm mem}(O)=
\left\{\rho\in\mathcal I_L:
\inf_{0\leq t\leq T_L}
\left|\Tr[O\rho(t)]-\Tr[O\rho_{\rm ref}]\right|\geq m_0
\right\}.
\]
\end{definition}

This records memory relative to a declared reference.  It does not by
itself establish MBL, exponential basin multiplicity, or a complete
LIOM algebra.

\section{Why there is no universal knee}

For finite $L$, \cref{eq:xxz} is a finite matrix depending analytically
on its parameters.  Away from degeneracies, its eigenvalues and spectral
projectors vary locally analytically.  Disorder-averaged finite-volume
diagnostics are therefore generally smooth functions of $W$ under mild
domination assumptions.  A sharp thermodynamic transition, if present,
appears only after a specified order of large-$L$, long-time, and
disorder limits.

The previous paper imported a Kramers escape law.  Kramers theory
requires a stochastic process, a barrier, a noise scale, and a
small-noise asymptotic.  None is supplied by the closed unitary
Hamiltonian \cref{eq:xxz}.  Coarse-graining may motivate a stochastic
effective model, but then its generator and approximation error must be
derived or declared.  A logistic curve is likewise only a fitting
family.  Under a monotone reparameterization of $W$, its midpoint and
width change, so they cannot be universal geometric invariants.

\begin{proposition}[Finite data do not select a unique crossover law]
\label{prop:interpolation}
Given finitely many distinct disorder values $W_1,\ldots,W_N$ and
diagnostic means $y_1,\ldots,y_N$, infinitely many smooth functions
$g(W)$ interpolate the data exactly while having different inflection
points and asymptotic behavior.
\end{proposition}

\begin{proof}
Let $p$ be the interpolating polynomial.  For any smooth $h$,
\[
g_h(W)=p(W)+h(W)\prod_{j=1}^{N}(W-W_j)
\]
matches every data point.  Varying $h$ changes the shape between and
outside the sampled values.
\end{proof}

The correct numerical question is not whether a curve looks like a
knee.  It is whether a declared scaling family gives stable,
cross-validated parameters as sizes and fit windows change.

\section{Finite-size and uncertainty protocol}

A reproducible comparison for \cref{eq:xxz} should report:
\begin{enumerate}
\item $L$, $J$, $\Delta$, boundary conditions, symmetry sector, energy
window, disorder distribution, random seeds, and sample counts;
\item $\langle r\rangle$, diagonal ETH residuals, imbalance, and
entanglement with sample-level data or sufficient statistics;
\item numerical residuals for eigenpairs and convergence checks;
\item disorder confidence intervals obtained from independent samples;
\item at least two plausible finite-size scaling models, including a
smooth-crossover null model;
\item sensitivity to omitted sizes, shifted energy windows, and changed
fit ranges; and
\item a held-out prediction, such as one larger size or a second
observable not used to locate the crossover.
\end{enumerate}

For a candidate dimensionless diagnostic $X_L(W)$, one may test
\[
 X_L(W)=F\!\left((W-W_c)L^{1/\nu}\right)
 +L^{-\omega}G\!\left((W-W_c)L^{1/\nu}\right),
\]
but this is a model, not a theorem.  The functions or their
parameterizations, the correction exponent, covariance model, and fit
window all belong to the parameter ledger.  Drift in fitted $W_c$ with
$L$ is evidence about finite-size effects, not permission to fix the
desired answer.  Large exact-diagonalization studies illustrate both the
power and limitations of this procedure \cite{luitz2015}.

\section{Interface with Modal Triplet Theory}

\subsection{What MTT supplies at its current tier}

The MTT foundation defines a coherent sector through spectral projection
and distinguishes stabilization, physical time, and coarse effective
dynamics \cite{nero_foundation}.  Its controlled-truncation paper states
the domain and resolvent data required for a valid coherent-sector
reduction \cite{nero_truncation}.  These structures motivate asking
whether a many-body experiment sees only a stable projected sector.
They do not, by themselves, select \cref{eq:xxz} or prove ETH or MBL.

In particular, the following gaps must not be identified by notation:
\begin{center}
\begin{tabular}{
>{\raggedright\arraybackslash}p{0.27\linewidth}
>{\raggedright\arraybackslash}p{0.29\linewidth}
>{\raggedright\arraybackslash}p{0.31\linewidth}}
\toprule
Object & Meaning & Required comparison \\
\midrule
MTT Hessian gap & local stability of a selected fixed-point problem &
an explicit source map and transported operator norm \\
Hamiltonian level spacing & differences among many-body energies &
symmetry-resolved spectral statistics \\
Liouvillian or channel gap & mixing rate of an open or coarse dynamics &
a declared generator or channel \\
mobility edge & energy-dependent localization boundary &
a thermodynamic diagnostic and scaling theorem \\
\bottomrule
\end{tabular}
\end{center}

\subsection{The missing pullback theorem}

Let $Q_{\rm MTT}$ denote a selected MTT coherent projector on an upper
space and let $\mathfrak C_L$ be the capsule in
\cref{def:capsule}.  A genuine MTT derivation would require a map
$V_L$ and channel-level construction satisfying, with controlled errors,
\begin{align}
 V_L Q_{\rm MTT} V_L^\dagger &\longmapsto
 \text{the selected effective sector},\\
 V_L A_{\rm MTT} V_L^\dagger &\longmapsto H_L(\omega;W),\\
 (\Pi_{\rm MTT},\text{margin}) &\longmapsto
 (\EE,\II,\Obs_L,\mathcal I_L,\mathcal D_L).
\end{align}
It must also preserve or quantitatively transport connections, domains,
holonomies, and the relevant operator bounds.  Only then could an MTT
stability statement imply one of the operational criteria proved here.

\begin{theorem}[Conditional diagnostic transport]\label{thm:transport}
Suppose a selected MTT model supplies maps into
$\mathfrak C_L$ such that:
\begin{enumerate}
\item its physical evolution is intertwined with $\Dyn_t$ up to channel
error $\epsilon_t$;
\item its selected coarse map is intertwined with $\EE$ and $\II$ up to
diamond-norm error $\epsilon_{\rm c}$; and
\item its selected observable $O^{\rm MTT}$ is transported to
$O\in\Obs_L$ with operator-norm error $\epsilon_O$.
\end{enumerate}
Then every finite-time expectation-value statement transports with an
explicit error bounded by the sum of the channel and observable errors.
In particular, for normalized states,
\[
 |\Delta\langle O\rangle|
 \leq \epsilon_t+\epsilon_{\rm c}+2\epsilon_O
\]
under a convention in which channel errors are evaluated in diamond
norm and $\norm{O}\leq1$.
\end{theorem}

\begin{proof}
Insert and subtract the exactly transported state and observable.  Use
trace duality for the observable terms and the defining diamond-norm
bound for the channel terms.  The two observable replacements contribute
at most $2\epsilon_O$; the evolution and coarse-channel replacements
contribute at most $\epsilon_t+\epsilon_{\rm c}$.
\end{proof}

This theorem says what a bridge can legitimately do.  It transports a
verified result after the source and error bounds are supplied.  It does
not create those source rows.

\section{Relation to monitored transitions}

Measurement-induced entanglement transitions involve stochastic
measurement records, conditional trajectories, nonlinear normalization,
and an ensemble prescription.  Their area- and volume-law regimes can
also be summarized by operational basins, but that shared vocabulary is
not a theorem that they are the same transition as ETH--MBL.  The
Hamiltonian channel \cref{eq:reduced} and a trajectory ensemble have
different state spaces and generators.  They should be compared only
after a separate typed map identifies the observables and limiting
procedures.  The companion paper on measurement-induced transitions is
the proper place for that construction.

\section{Falsifiers and next calculations}

The operational bridge is useful because it can fail cleanly.
\begin{enumerate}
\item If a proposed effective channel is claimed to be Markovian but
\cref{eq:defect} is not small on the tested states, the reduction fails.
\item If contraction is claimed to explain thermalization while
microscopic trace distance is said to contract under unitary evolution,
\cref{prop:isometry} identifies the type error.
\item If a thermal basin is claimed but either term in
\cref{eq:decomposition} remains large, the thermalization claim fails for
that state and observable.
\item If a quasi-LIOM is claimed, its commutator norm must support the
reported memory time through \cref{eq:liom}.
\item If a universal knee is claimed, stability under alternative scaling
models and reparameterizations must be demonstrated on held-out sizes.
\item If the MTT explanation is claimed, the source and intertwining rows
in \cref{thm:transport} must be emitted from selected geometry rather
than fitted to the many-body output.
\end{enumerate}

The immediate computational program is therefore modest: construct
$\mathfrak C_L$ for a sequence of tractable sizes, publish sample-level
diagnostics and uncertainty, quantify the semigroup defect for one
declared coarse channel, and test whether the resulting operational
basin labels add predictive information beyond the standard diagnostics.
No numerical result is claimed in this paper.

\section{Conclusion}

ETH and MBL can be discussed in one operational language, but they have
not been proved to be one universal projection phenomenon.  The useful
common structure is a declared record of states, observables, time
windows, coarse maps, and errors.  Within that record, uniform
contraction is a sufficient mixing criterion; diagonal bias and temporal
fluctuation are exactly separable; and an approximately conserved
quasi-local operator yields a quantitative memory time.  Those are
rigorous bridges because their assumptions and conclusions are visible.

MTT adds a potentially deeper question: can one selected upper geometry
emit the Hamiltonian sector, coarse channel, observables, and stability
data together?  At present that is an open pullback theorem.  Stating it
explicitly strengthens rather than weakens the program: it distinguishes
an explanatory source from an adaptable after-the-fact description and
turns the next step into a concrete mathematical construction.

\bibliographystyle{plain}
\bibliography{main}

\end{document}
